% =====================================================================
%  Lecture 25 -- Network Evolution: a Grand Synthesis    (Persian edition)
% =====================================================================
\documentclass[11pt,a4paper]{article}
\providecommand\ACNroot{../..}
\input{\ACNroot/templates/lecture-fa}

\ACNset{number}{25}
\ACNset{title}{تکامل شبکه: یک جمع‌بندی بزرگ}
\ACNset{shorttitle}{تکامل شبکه}
\ACNset{titleother}{Network Evolution: a Grand Synthesis}
\ACNset{part}{بخش پنجم -- صفحه‌ها و چشم‌انداز (درس‌های ۲۳ تا ۲۶)}
\ACNset{time}{۶ تا ۸ ساعت (۳ ساعت مطالعه، ۳ تا ۵ ساعت تمرین)}
\ACNset{level}{پیشرفته}
\ACNset{prereq}{درس‌های ۱ تا ۲۴}
\ACNset{gen}{5}

\begin{document}
\ACNcover

% ---------------------------------------------------------------------
\section{اهداف یادگیری}
پس از مطالعهٔ این درس باید بتوانید:
\begin{objectives}
  \item تکامل دسترسی رادیویی، معماری هسته، سیگنالینگ، امنیت و \acro{QoS} را از نسل دوم تا نسل پنجم \textbf{دنبال کنید}.
  \item کارکردهای شبکه را میان نسل‌ها \textbf{نگاشت کنید} (\lr{MSC/VLR}، \acro{MME}، \acro{AMF}؛ \acro{HLR}، \acro{HSS}، \acro{UDM}؛ \acro{GGSN}، \lr{P-GW}، \acro{UPF}).
  \item رشد ظرفیت را به طیف، بازدهی طیفی و متراکم‌سازی \textbf{تجزیه کنید}.
  \item روندهای بلندمدت نرخ اوج، تأخیر و هزینه به ازای بیت را با مدل‌های رشد ساده \textbf{کمّی کنید}.
  \item اصول طراحی‌ای را که در هر نسل تکرار می‌شوند و نسل بعد را شکل خواهند داد \textbf{شناسایی کنید}.
\end{objectives}

% ---------------------------------------------------------------------
\section{دانش پیش‌نیاز}
این درس از کل دوره بهره می‌گیرد: \acro{GSM}، \acro{GPRS} و هستهٔ نسل دوم و سوم (درس‌های ۱ تا ۴)، \acro{LTE} و \acro{EPC} (درس‌های ۵ تا ۱۳) و سامانهٔ
نسل پنجم، \acro{SDN}، \acro{NFV} و \acro{SBA} (درس‌های ۱۴ تا ۲۴). هیچ پروتکل جدیدی معرفی نمی‌شود؛ هدف پیوند دادن آنچه می‌دانید در یک تصویر واحد است.

\paragraph{واژه‌های کلیدی} واژگان این درس همراه با معادل انگلیسی آن‌ها:
\begin{keyterms}
نسل & \lr{Generation} & خانواده‌ای از استانداردها که مجموعه‌ای از الزامات \acro{ITU} را برآورده می‌کند (\lr{IMT-2000}، \lr{IMT-Advanced}، \lr{IMT-2020}). \\
معماری تخت & \lr{Flat architecture} & گره‌های کمتر در مسیر کاربر؛ کنترل رادیو در ایستگاه پایه. \\
تمام‌\lr{IP} & \lr{All-IP} & حمل صدا و داده به‌صورت بستهٔ \acro{IP}، بدون حوزهٔ سوئیچینگ مداری. \\
جداسازی صفحهٔ کنترل و کاربر & \lr{CP/UP separation} & مقیاس‌پذیری و جانمایی مستقل صفحهٔ کنترل و کاربر. \\
بازدهی طیفی & \lr{Spectral efficiency} & گذردهی به ازای هر هرتز در هر سلول (\lr{bit/s/Hz/cell}). \\
متراکم‌سازی & \lr{Densification} & سلول‌های بیشتر در واحد سطح. \\
ظرفیت ترافیکی ناحیه & \lr{Area traffic capacity} & گذردهی کل در واحد سطح (\lr{bit/s/m}$^2$). \\
هزینه به ازای بیت & \lr{Cost per bit} & هزینهٔ شبکه تقسیم بر ترافیک حمل‌شده. \\
نرم‌افزاری‌شدن & \lr{Softwarisation} & کارکردهای شبکه به‌صورت نرم‌افزار روی زیرساخت همه‌منظوره. \\
سازگاری با گذشته & \lr{Backward compatibility} & همکاری یک نسل جدید با نسل‌های نصب‌شده. \\
\end{keyterms}

% ---------------------------------------------------------------------
\section{مقدمه و انگیزه}
هر نسل از شبکه‌های موبایل برای رفع گلوگاه نسل پیشین طراحی شد: نسل دوم صدا را دیجیتال کرد و امنیت و رومینگ افزود؛ نسل سوم دادهٔ بسته‌ای با
نرخ‌های کاربردی آورد؛ نسل چهارم حوزهٔ مداری را حذف و شبکه را تخت کرد؛ نسل پنجم شبکه را برنامه‌پذیر و خدمت‌محور کرد. اگر نسل‌ها را یکی‌یکی
ببینیم، مانند دنباله‌ای از پشته‌های پروتکلی بی‌ارتباط به نظر می‌رسند. اگر با هم ببینیم، مجموعهٔ کوچکی از ایده‌های پایدار را نشان می‌دهند --
جداسازی دغدغه‌ها، لنگرهای جابه‌جایی، پایگاه‌های دادهٔ مشترک، تونل‌زنی، احراز هویت متقابل -- که با فناوری‌هایی هر بار عمومی‌تر پیاده شده‌اند
\cite{itu2015m2083,ts23501}. شناختن این ایده‌هاست که به یک مهندس اجازه می‌دهد سامانهٔ جدید را سریع بفهمد و مسیر آینده‌اش را پیش‌بینی کند.

% ---------------------------------------------------------------------
\section{مفاهیم بنیادی}

\subsection{پنج بُعد تکامل}
\begin{itemize}
  \item \textbf{دسترسی رادیویی}: \lr{FDMA/TDMA} (\acro{GSM})، \acro{CDMA} (\acro{UMTS})، \acro{OFDMA} با \acro{MIMO} (\acro{LTE})، \acro{OFDM} با عددشناسی
        انعطاف‌پذیر همراه با \acro{MIMO} انبوه و موج میلی‌متری (\acro{NR}).
  \item \textbf{هسته}: سوئیچینگ مداری با \lr{MSC/VLR} و \lr{SS7/MAP}؛ یک هستهٔ بسته‌ای روکش‌شده (\lr{SGSN/GGSN})؛ \acro{EPC} تمام‌\acro{IP} با
        \lr{Diameter} و \acro{GTP}؛ \lr{5GC} خدمت‌محور روی زیرساخت ابری.
  \item \textbf{جابه‌جایی}: نواحی مکان و مسیریابی سلسله‌مراتبی، لنگرها در \acro{GGSN}، \lr{P-GW} و \acro{UPF}؛ تصمیم‌های دست‌به‌دست که از \acro{BSC}
        و \acro{RNC} به ایستگاه پایه منتقل شدند.
  \item \textbf{امنیت}: احراز هویت یک‌طرفهٔ \acro{GSM} (\lr{A3/A8}) تا \acro{AKA} متقابل \acro{UMTS}، \lr{EPS-AKA} با سلسله‌مراتب کلید، و
        \lr{5G-AKA}/\lr{EAP-AKA$'$} با شناسه‌های پنهان (\acro{SUCI}) و کنترل شبکهٔ خانگی.
  \item \textbf{\acro{QoS}}: از تضمین‌های مداری از طریق زمینه‌های \acro{PDP} و حامل‌های \acro{EPS} تا \acro{QoS} جریان‌محور با \acro{5QI}های
        استاندارد و برش‌بندی.
\end{itemize}

\subsection{تخت‌شدن مسیر}
مسیر صفحهٔ کاربر پیوسته کوتاه‌تر شده است (\cref{fig:arch}). در \acro{GPRS} و \acro{UMTS} یک بسته از یک ایستگاه پایه، یک کنترل‌گر رادیو
(\lr{BSC/PCU} یا \acro{RNC})، یک \acro{SGSN} و یک \acro{GGSN} می‌گذشت. \acro{LTE} کنترل رادیو را به \acro{eNB} برد و کنترل‌گر را حذف کرد؛ \lr{S-GW} و
\lr{P-GW} باقی ماندند. در نسل پنجم یک \acro{UPF} می‌تواند کل صفحهٔ کاربر هسته باشد و نزدیک رادیو قرار گیرد.

\begin{acnfigure}
  \resizebox{\linewidth}{!}{\acnfig{lecture-25/arch-evolution}}
  \caption{مسیر صفحهٔ کاربر دادهٔ بسته‌ای از نسل ۲٫۵ تا نسل پنجم (ساده‌شده).}
  \label{fig:arch}
\end{acnfigure}

\subsection{کارکردهایی که می‌مانند}
نام‌ها تغییر می‌کنند، کارکردها می‌مانند. کارکرد مکان بازدیدکننده و جابه‌جایی (\lr{MSC/VLR}، \acro{SGSN}، \acro{MME}، \acro{AMF})، پایگاه دادهٔ مشترک
(\lr{HLR/AuC}، \acro{HSS}، \lr{UDM/AUSF})، لنگر جابه‌جایی و دروازه به شبکه‌های داده (\acro{GGSN}، \lr{P-GW}، \acro{UPF} همراه با \acro{SMF}) و کنترل سیاست
(\acro{PCRF}، \acro{PCF}) در همهٔ نسل‌ها حضور دارند، هر بار با واسط‌های تمیزتر و جداسازی بهتر صفحهٔ کنترل و کاربر.

% ---------------------------------------------------------------------
\section{توضیحات فنی تفصیلی}

\subsection{از مدار به بسته به خدمت}
نسل دوم حول یک مدار ۶۴ کیلوبیت بر ثانیه برای هر تماس طراحی شد. \acro{GPRS} یک هستهٔ بسته‌ای موازی افزود و رادیو را با تخصیص پویای
شکاف زمانی دوباره به کار گرفت. \acro{UMTS} دو حوزه (\acro{CS} و \acro{PS}) را با یک شبکهٔ رادیویی مشترک نگه داشت. \acro{LTE} حوزهٔ \acro{CS} را حذف
کرد: صدا به یک خدمت \acro{IMS} روی \acro{VoLTE} تبدیل شد، با \lr{CS fallback} و \acro{SRVCC} برای همکاری. نسل پنجم \acro{IMS} را برای صدا (\acro{VoNR})
نگه می‌دارد اما هسته را به‌صورت خدماتی که روی \lr{HTTP/2} مصرف می‌شوند بازسازی می‌کند -- خود شبکه به مجموعه‌ای از \acro{API}ها تبدیل می‌شود.

\subsection{پروتکل‌های سیگنالینگ}
\acro{SS7} و \acro{MAP} هستهٔ نسل دوم و سوم را به هم وصل می‌کردند؛ \lr{Diameter} در \acro{EPC} جای \acro{MAP} را گرفت (\lr{S6a}، \lr{Gx}، \lr{Rx})؛
\lr{GTPv2-C} کنترل نشست را حمل می‌کرد؛ نسل پنجم از واسط‌های خدمت‌محور به سبک \acro{REST} با مجوزدهی مبتنی بر \lr{OAuth} و از \acro{PFCP} برای
کنترل صفحهٔ کاربر استفاده می‌کند. هر گام به سوی پروتکل‌هایی رفت که در صنعت گسترده‌تر فناوری اطلاعات به کار می‌روند، و این هزینهٔ توسعه را کاهش
داد و استقرار ابری را ممکن کرد.

\subsection{درس‌های امنیتی}
\acro{GSM} تنها مشترک را احراز هویت می‌کرد و این امکان ایستگاه پایهٔ جعلی را می‌داد؛ \acro{UMTS} احراز هویت شبکه و حفاظت یکپارچگی را افزود؛
\acro{LTE} کلیدها را برای هر لایه و هر گره جدا کرد؛ نسل پنجم شناسهٔ دائمی را رمز می‌کند (\acro{SUCI})، تصمیم نهایی احراز هویت را به شبکهٔ خانگی
می‌دهد و سیگنالینگ میان اپراتورها را با \acro{SEPP} حفاظت می‌کند. هر اصلاح پاسخی به حمله‌ای بود که در نسل پیشین عملی شده بود.

\subsection{هم‌زیستی}
هیچ نسلی نسل پیش از خود را یک‌باره جایگزین نمی‌کند. شبکه‌ها یک دهه یا بیشتر چند نسل را به‌موازات اجرا می‌کنند، با کارکردهای همکاری (\lr{N26} میان
\acro{AMF} و \acro{MME}، \lr{EN-DC} در نسل پنجم غیرمستقل، \lr{CS fallback}) و بازتخصیص تدریجی طیف. سازگاری با گذشته یک قید طراحی همیشگی و بخش
بزرگی از هزینهٔ هر سامانهٔ جدید است.

% ---------------------------------------------------------------------
\section{مدل‌های ریاضی}

\subsection{رشد نرخ اوج}\label{sec:peak}
\paragraph{مدل} یک روند نمایی $R(t) = R_0\, e^{k (t - t_0)}$ که از دو نقطهٔ $(t_0, R_0)$ و $(t_1, R_1)$ برازش شود دارای
\begin{equation}\label{eq:growth}
  k = \frac{\ln (R_1/R_0)}{t_1 - t_0}, \qquad T_2 = \frac{\ln 2}{k}
\end{equation}
است.
\paragraph{مثال} از دادهٔ \acro{GSM} (۹٫۶ کیلوبیت بر ثانیه، ۱۹۹۱) تا الزام \lr{IMT-2020} (۲۰ گیگابیت بر ثانیه، تدوین ۲۰۱۷ تا ۲۰۱۸) \cite{itu2017m2410}:
$k = \ln(2.08\times10^{6})/27 = 0.539$ در سال، یعنی ضریب ۱٫۷۱ در سال و دو برابر شدن هر ۱٫۲۹ سال (\cref{fig:peak}).
\paragraph{محدودیت} نرخ‌های اوج مقادیر مشخصات در شرایط ایده‌آل‌اند؛ نرخ تجربه‌شدهٔ کاربر کندتر رشد می‌کند و به بار و پوشش بستگی دارد.

\begin{acnfigure}
  \acnfig{lecture-25/peak-rates}
  \caption{نرخ‌های اوج فروسوی تعیین‌شدهٔ فناوری‌های \lr{3GPP} و برازش نمایی \cref{eq:growth}.}
  \label{fig:peak}
\end{acnfigure}

\subsection{تجزیهٔ ظرفیت}\label{sec:cap}
\paragraph{مدل} ظرفیت ترافیکی ناحیه حاصل‌ضرب چگالی سلول $n$ (سلول بر کیلومتر مربع)، پهنای باند $B$ و بازدهی طیفی $\eta$ است:
\begin{equation}\label{eq:area}
  C_A = n\, B\, \eta ,
\end{equation}
پس رشد میان دو استقرار حاصل‌ضرب سه نسبت است.
\paragraph{مثال} پایهٔ \acro{LTE}: $n = 4$، $B = 20$ مگاهرتز، $\eta = 1.5$: $C_A = 120$ مگابیت بر ثانیه بر کیلومتر مربع. استقرار نسل پنجم: $n = 16$،
$B = 100$ مگاهرتز، $\eta = 6$: $C_A = 9.6$ گیگابیت بر ثانیه بر کیلومتر مربع -- ضریب $4\times5\times4 = 80$ (\cref{fig:cap}). برای مقایسه، \lr{IMT-2020} در سناریوی
نقطهٔ داغ درون‌ساختمانی ۱۰ مگابیت بر ثانیه بر متر مربع (۱۰ ترابیت بر ثانیه بر کیلومتر مربع) را الزامی می‌کند \cite{itu2017m2410}.

\begin{acnfigure}
  \acnfig{lecture-25/capacity-decomp}
  \caption{رشد ظرفیت ناحیه از پایهٔ \acro{LTE} تا استقرار نسل پنجم در \cref{sec:cap}.}
  \label{fig:cap}
\end{acnfigure}

\subsection{تأخیر در طول مسیر}\label{sec:path}
\paragraph{مدل} تأخیر یک‌طرفهٔ صفحهٔ کاربر از میان $h$ گره برابر است با
\begin{equation}\label{eq:path}
  T = \sum_{i=1}^{h} \left( t_{\mathrm{proc},i} + \frac{d_i}{v} \right) ,
\end{equation}
با زمان پردازش $t_{\mathrm{proc},i}$، طول پیوند $d_i$ و سرعت انتشار $v \approx 2\times10^{8}$ متر بر ثانیه در فیبر.
\paragraph{مثال} چهار گره هر کدام ۰٫۵ میلی‌ثانیه و ۴۰۰ کیلومتر فیبر (۲ میلی‌ثانیه) ۴ میلی‌ثانیه می‌دهند؛ دو گره (\acro{gNB} و \acro{UPF} لبه) هر کدام ۰٫۲۵
میلی‌ثانیه و ۲۰ کیلومتر ۰٫۶ میلی‌ثانیه می‌دهند. تخت‌سازی و جانمایی لبه با هم بیشتر تأخیر هسته را حذف می‌کنند.

\subsection{هزینه به ازای بیت}\label{sec:cost}
\paragraph{مدل} اگر ترافیک سالانه با ضریب $1+g$ و درآمد با ضریب $1+r$ رشد کند، هزینه به ازای بیت باید سالانه با ضریب $(1+r)/(1+g)$ کاهش یابد تا حاشیهٔ سود
ثابت بماند؛ این هزینه هر
\begin{equation}\label{eq:cost}
  T_{1/2} = \frac{\ln 2}{\ln (1+g) - \ln (1+r)} \ \text{\rl{سال}}
\end{equation}
نصف می‌شود.
\paragraph{مثال} با $g = 30\,\%$ و درآمد ثابت، $T_{1/2} = 2.64$ سال. این فشار، بیش از نرخ اوج، حرکت صنعت به سوی سخت‌افزار عمومی، خودکارسازی و
شبکه‌های نرم‌افزارمحور را توضیح می‌دهد.

% ---------------------------------------------------------------------
\section{معماری شبکه و نمودارها}
\cref{fig:arch} مسیر کاربر را نشان می‌دهد؛ سوی کنترل به‌موازات از \lr{MSC/VLR} و \acro{SGSN} به \acro{MME} و سپس به \acro{AMF} و \acro{SMF} در هستهٔ
خدمت‌محور تکامل یافت. \cref{tab:map} کارکردها را میان نسل‌ها نگاشت می‌کند.

\begin{table}[htbp]
\centering\small
\caption{کارکردهای متناظر در نسل‌های مختلف.}\label{tab:map}
\rowcolors{2}{acnmist}{white}
\begin{tabularx}{\linewidth}{@{}l Y Y Y@{}}
\toprule
\rowcolor{white} کارکرد & \lr{2G/3G} & \lr{4G (EPC)} & \lr{5G (5GC)} \\
\midrule
جابه‌جایی و کنترل دسترسی & \lr{MSC/VLR}، \acro{SGSN} & \acro{MME} & \acro{AMF} \\
دادهٔ مشترک و احراز هویت & \lr{HLR/AuC} & \acro{HSS} & \acro{UDM}، \acro{AUSF} \\
کنترل نشست & \acro{SGSN}، \acro{GGSN} & \acro{MME}، کنترل \lr{S-GW/P-GW} & \acro{SMF} \\
لنگر صفحهٔ کاربر & \acro{GGSN} & \lr{P-GW} (\lr{S-GW}) & \acro{UPF} \\
سیاست & -- (ایستا) & \acro{PCRF} & \acro{PCF} \\
کنترل رادیو & \acro{BSC}، \acro{RNC} & \acro{eNB} & \acro{gNB} (\lr{CU-CP}) \\
\bottomrule
\end{tabularx}
\end{table}

% ---------------------------------------------------------------------
\section{رویه‌های گام‌به‌گام}

\begin{procedurebox}{تحلیل یک سامانهٔ موبایل جدید}\label{proc:analyse}
\begin{steps}
  \item پایگاه دادهٔ مشترک، کارکرد جابه‌جایی، کارکرد نشست و لنگر صفحهٔ کاربر را شناسایی کنید.
  \item مسیر کاربر را رسم و گره‌ها و تونل‌ها را بشمارید (\cref{eq:path}).
  \item روش احراز هویت و شناسه‌هایی را که به‌صورت آشکار منتقل می‌شوند شناسایی کنید.
  \item واحد \acro{QoS} (مدار، حامل، جریان) و محل اعمال آن را شناسایی کنید.
  \item همکاری با نسل پیشین را پیدا کنید.
\end{steps}
\end{procedurebox}

\begin{procedurebox}{برنامه‌ریزی افزایش ظرفیت}\label{proc:cap}
\begin{steps}
  \item ترافیک کنونی را اندازه‌گیری و رشد را پیش‌بینی کنید.
  \item ضریب لازم را محاسبه و آن را میان طیف، بازدهی و چگالی تقسیم کنید (\cref{eq:area}).
  \item هزینهٔ هر اهرم را برآورد کنید (پروانه‌ها، تجهیزات جدید، سایت‌ها و بک‌هال).
  \item ارزان‌ترین ترکیب را برگزینید و هدف هزینه به ازای بیت را بررسی کنید (\cref{eq:cost}).
\end{steps}
\end{procedurebox}

% ---------------------------------------------------------------------
\section{مثال‌های حل‌شده}

\begin{examplebox}{متراکم‌سازی لازم}\label{ex:dens}
\textbf{مسئله.} یک نقطهٔ داغ باید ۱۰ ترابیت بر ثانیه بر کیلومتر مربع را با $B = 400$ مگاهرتز و $\eta = 10$ حمل کند. چند سلول در هر کیلومتر مربع لازم است؟

\textbf{حل.} $n = 10^{13}/(4\times10^{8}\times10) = 2500$ سلول بر کیلومتر مربع، یک سلول برای هر ۴۰۰ متر مربع -- فاصلهٔ میان سایت‌ها حدود ۲۰ متر، یعنی سلول‌های
کوچک درون‌ساختمانی.
\end{examplebox}

\begin{examplebox}{برون‌یابی روند}\label{ex:extra}
\textbf{مسئله.} اگر برازش \cref{sec:peak} ادامه می‌یافت، نرخ اوج کی به ۱ ترابیت بر ثانیه -- رقمی که در چشم‌اندازهای نسل ششم آمده \cite{saad2020vision} -- می‌رسید؟

\textbf{حل.} $\ln(10^{12}/2\times10^{10})/0.539 = 7.3$ سال پس از ۲۰۱۸، یعنی حدود ۲۰۲۵. چنین نشده است، که نشان می‌دهد این روند قانون نیست: رشد بیشتر به پهنای
باندهای گسترده‌تری نیاز دارد که تنها بالای ۱۰۰ گیگاهرتز وجود دارند (درس ۲۶).
\end{examplebox}

\begin{examplebox}{فشار هزینه}\label{ex:cost}
\textbf{مسئله.} ترافیک سالانه ۲۵٪ و درآمد ۲٪ رشد می‌کند. هزینه به ازای بیت هر سال چقدر باید کاهش یابد و کی نصف می‌شود؟

\textbf{حل.} $1.02/1.25 = 0.816$: کاهش سالانهٔ ۱۸٫۴٪؛ $T_{1/2} = \ln 2/(\ln 1.25 - \ln 1.02) = 3.41$ سال.
\end{examplebox}

% ---------------------------------------------------------------------
\section{مقایسهٔ فناوری‌ها}
\begin{table}[htbp]
\centering\small
\caption{نسل‌ها در یک نگاه.}\label{tab:cmp}
\rowcolors{2}{acnmist}{white}
\begin{tabularx}{\linewidth}{@{}l Y Y Y Y@{}}
\toprule
\rowcolor{white} & \lr{2G (GSM)} & \lr{3G (UMTS)} & \lr{4G (LTE)} & \lr{5G (NR)} \\
\midrule
دسترسی چندگانه & \lr{TDMA/FDMA} & \acro{WCDMA} & \lr{OFDMA / SC-FDMA} & \acro{OFDMA}، عددشناسی انعطاف‌پذیر \\
هسته & \acro{CS} (+ \acro{GPRS}) & \lr{CS + PS} & \acro{EPC} تمام‌\acro{IP} & \lr{5GC} خدمت‌محور \\
سیگنالینگ & \lr{SS7/MAP} & \lr{SS7/MAP} & \lr{Diameter}، \lr{GTPv2-C} & \acro{SBI} روی \lr{HTTP/2}، \acro{PFCP} \\
احراز هویت & یک‌طرفه & \acro{AKA} متقابل & \lr{EPS-AKA} & \lr{5G-AKA}، \lr{EAP-AKA$'$}، \acro{SUCI} \\
واحد \acro{QoS} & مدار & زمینهٔ \acro{PDP} & حامل \acro{EPS} & جریان \acro{QoS}، برش \\
صدا & \acro{CS} & \acro{CS} & \lr{VoLTE (IMS)} & \lr{VoNR (IMS)} \\
\bottomrule
\end{tabularx}
\end{table}

\paragraph{همان ایده‌ها، ابزارهای بهتر} هر نسل همان نقش‌های کارکردی را نگه می‌دارد اما آن‌ها را با فناوری عمومی‌تر بازمی‌سازد: از سوئیچ‌های اختصاصی
به مسیریاب‌های \acro{IP} و سپس به نرم‌افزار ابری.

% ---------------------------------------------------------------------
\section{نکات کلیدی}
\begin{takeawaybox}
\begin{itemize}
  \item کارکردها میان نسل‌ها پایدارند؛ نام، پروتکل و جانمایی آن‌ها تغییر می‌کند.
  \item مسیر کاربر از چهار گره (\lr{GPRS/UMTS}) به یک \acro{UPF} در نسل پنجم تخت شد.
  \item ظرفیت $C_A = nB\eta$: رشد از طیف، بازدهی و متراکم‌سازی با هم می‌آید.
  \item نرخ اوج از ۱۹۹۱ تا ۲۰۱۸ حدود ۱٫۷ برابر در سال رشد کرد، اما روندها قانون نیستند.
  \item هزینه به ازای بیت باید هر $\ln 2/(\ln(1+g) - \ln(1+r))$ سال نصف شود -- محرک نرم‌افزاری‌شدن.
\end{itemize}
\end{takeawaybox}

% ---------------------------------------------------------------------
\section{تمرین‌ها}

\begin{exercisebox}{1}{نگاشت}\label{exr:map}
معادل نسل چهارم و پنجم این‌ها را نام ببرید: (الف) \acro{VLR}؛ (ب) \acro{HLR}؛ (ج) \acro{GGSN}؛ (د) \acro{RNC}.
\end{exercisebox}

\begin{exercisebox}{1}{چرا بسته؟}\label{exr:why}
دو دلیل بیاورید که چرا شبکه‌های موبایل برای همهٔ ترافیک از سوئیچینگ مداری به بسته‌ای رفتند.
\end{exercisebox}

\begin{exercisebox}{2}{نرخ رشد}\label{exr:cagr}
\lr{LTE Rel-8} ۳۰۰ مگابیت بر ثانیه (۲۰۰۸) و \lr{LTE-Advanced} ۳ گیگابیت بر ثانیه (۲۰۱۱) ارائه کردند. $k$ و زمان دو برابر شدن را برای این بازه محاسبه و با
\cref{sec:peak} مقایسه کنید.
\end{exercisebox}

\begin{exercisebox}{2}{اهرم‌های ظرفیت}\label{exr:cap}
یک اپراتور به ۲۰ برابر ظرفیت ناحیه نیاز دارد. طیف می‌تواند ۲٫۵ برابر و بازدهی ۲ برابر شود. چه میزان متراکم‌سازی لازم است؟
\end{exercisebox}

\begin{exercisebox}{3}{تأخیر مسیر}\label{exr:path}
با \cref{eq:path} مقایسه کنید: (الف) مسیر \acro{UMTS} با \lr{NodeB}، \acro{RNC}، \acro{SGSN} و \acro{GGSN} (هر کدام ۱ میلی‌ثانیه، ۳۰۰ کیلومتر فیبر)؛ (ب) نسل
پنجم با \acro{gNB} و یک \acro{UPF} در ۳۰ کیلومتری (هر کدام ۰٫۳ میلی‌ثانیه).
\end{exercisebox}

\begin{exercisebox}{3}{گلوگاه بعدی}\label{exr:next}
با استفاده از پنج بُعد تکامل، با یک استدلال کمّی بحث کنید نسل ششم به احتمال زیاد کدام گلوگاه نسل پنجم را هدف می‌گیرد.
\end{exercisebox}

\subsection*{پاسخ‌ها}
\addcontentsline{toc}{subsection}{پاسخ‌ها}

\begin{solutionbox}{\ref{exr:map}}
(الف) \acro{MME} / \acro{AMF}؛ (ب) \acro{HSS} / \acro{UDM} (با \acro{AUSF})؛ (ج) \lr{P-GW} / \acro{UPF} همراه با \acro{SMF}؛ (د) گرهٔ جداگانه‌ای ندارد --
\acro{eNB} / \acro{gNB} (\lr{CU-CP}).
\end{solutionbox}

\begin{solutionbox}{\ref{exr:why}}
تسهیم آماری ظرفیت را برای ترافیک انفجاری بسیار بهتر به کار می‌گیرد، و یک شبکهٔ \acro{IP} واحد برای صدا و داده هزینه را کاهش می‌دهد و ارائهٔ مستقیم
خدمات اینترنتی را ممکن می‌کند.
\end{solutionbox}

\begin{solutionbox}{\ref{exr:cagr}}
$k = \ln 10/3 = 0.768$ در سال، $T_2 = 0.90$ سال -- سریع‌تر از ۱٫۲۹ سال بلندمدت، چون تجمیع حامل و \acro{MIMO} پهنای باند و لایه‌ها را یک‌باره افزودند.
\end{solutionbox}

\begin{solutionbox}{\ref{exr:cap}}
$20/(2.5\times2) = 4$: چهار برابر سلول در واحد سطح.
\end{solutionbox}

\begin{solutionbox}{\ref{exr:path}}
(الف) $4\times1 + 300\,\text{km}/(2\times10^{5}\,\text{km/s}) = 4 + 1.5 = 5.5$ میلی‌ثانیه. (ب) $2\times0.3 + 0.15 = 0.75$ میلی‌ثانیه.
\end{solutionbox}

\begin{solutionbox}{\ref{exr:next}}
یک استدلال نمونه: روند نرخ اوج متوقف شد چون پهنای باند زیر ۱۰۰ گیگاهرتز تمام شده است؛ $C = B\log_2(1+\mathrm{SNR})$ با $B$ خطی ولی با \acro{SNR} تنها
لگاریتمی رشد می‌کند، پس گام ده‌برابری بعدی به ده‌ها گیگاهرتز طیف (زیرتراهرتز)، به‌علاوهٔ راه‌های جدید کاهش هزینه به ازای بیت (خودکارسازی مبتنی بر
هوش مصنوعی) و پوشش مناطق بدون اتصال (شبکه‌های غیرزمینی) نیاز دارد.
\end{solutionbox}

% ---------------------------------------------------------------------
\section{مطالعهٔ بیشتر}
\lr{ITU-R M.2083} \cite{itu2015m2083} و \lr{M.2410} \cite{itu2017m2410} چشم‌انداز و الزامات \lr{IMT-2020} را تعریف می‌کنند؛ \lr{TS 23.002} \cite{ts23002}،
\lr{TS 23.401} \cite{ts23401} و \lr{TS 23.501} \cite{ts23501} معماری نسل‌های پیاپی را شرح می‌دهند. ابرشپخر و همکاران \cite{eberspaecher2009gsm}، هولما و
توسکالا \cite{holma2010wcdma}، سزیا و همکاران \cite{sesia2011lte} و دالمن و همکاران \cite{dahlman2020nr} به \acro{GSM}، \acro{UMTS}، \acro{LTE} و \acro{NR}
می‌پردازند؛ رومر و همکاران \cite{rommer2020core} به هستهٔ نسل پنجم.

\begin{latin}
\printbibliography[heading=subbibliography,title={References}]
\end{latin}
\end{document}
